\chapter{Don't forget (Puella Magi Kazumi Magica)}
There are three major tenets established over the course of the Madoka Magica television series that, together, pose something of a challenge for each of the three spinoffs I'll be addressing in this section. First, magical girls are and always have been witches, dating back to the origin of magical girls. Given the original nature of witches as an image of the patriarchal fear of feminine power \cite{ref103}, this means magical girls within Madoka Magica are both inherently a challenge to the dominant power structure and easily readable as a corrupt and corrupting element. Second, the world of Madoka Magica is dominated by the Incubators' system, in which cosmic entropy is cleansed by converting it to the emotional decay of magical girls into witches without their consent—that is, the dominant power structure is exploitative of and deeply unjust toward women. Third, this power structure is highly stable, and is not overthrown until Madoka makes her wish and creates her new system.

Leaving for now the question of whether Madoka's system is any more just than Kyubey's (but see Interlude 2 and Chapter 17), the third tenet means that any spinoff work set during the show's main timeline or any of the "prior" timelines in which Kyubey's system dominates must ultimately treat Kyubey's system as an inescapable and unchangeable fact. Add in the ways in which Kyubey's system constrains and exploits the power of the magical girls (the second tenet) and the equation of that power to a tradition that depicts feminine power as corrupt and dangerous (the first tenet), and it becomes quite easy for a spinoff to fall into the trap of depicting the magical girls as dangerous and Kyubey's system as a necessary evil to contain them, or worse, a net positive.

This is precisely the trap into which the manga series Puella Magi Kazumi Magica (written by Masaki Hiramatsu, art by Takashi Tensugi) falls.

Kazumi Magica opens with the titular Kazumi as a near-total blank slate, naked, sealed in a trunk, with no memory of her past. This is an extreme form of what the infamous Turkey City Lexicon \cite{ref104} refers to as "White Room Syndrome," in which a story opens with a lone character in a white room wondering how they came to be there; the Lexicon rather uncharitably, but perhaps accurately, makes the assumption that this is because the author has yet to figure out the setting or characters. It is unclear whether this critique applies to Kazumi Magica; on the one hand, Kazumi's lack of memory serves an important plot function, since it allows her to serve as a vehicle for the audience to discover the truth about her identity and the Pleiades magical girl team. However, Kazumi Magica was a fairly long-running series at 23 monthly chapters, and began well before the series ended. Many of the later revelations in Kazumi Magica logically depend on late-series revelation in Madoka Magica, and it is unclear whether these revelations were held back to avoid the spinoff comic spoiling reveals in its parent series, or because Hiramatsu and Tensegu were unaware of the relevant spoilers and had to alter their comic's course in mid-stream to avoid contradicting the parent series. 

From here the comic moves into the reveal that Kazumi and her friends are magical girls, members of a seven-member team called the Pleiades, and supported by a seemingly friendlier variant of Kyubey named Juubey. As the comic unfolds, the reader learns that these Pleiades are far darker than they initially appear. Ultimately, it is revealed that they were originally a cooperative group of magical girls for whom Kazumi served as a mentor, but she eventually transformed into a witch. They combined their powers to transform Kazumi back into a humanlike form, and also killed and converted their local instance of Kyubey into Juubey, at the same time covering the entire city in a spell that renders the replacement Kyubey imperceptible. The Pleiades erroneously believe that Juubey can clean their soul gems for them, but in fact he only creates the illusion of cleansing; underneath, they are still slowly descending into witchery. Kazumi, meanwhile, is only the latest in a series of attempts to resurrect the original; each one eventually breaks down and becomes a witch. In the finale, all but two of the Pleiades are killed and a massive conglomerate witch created, but Kazumi comes to realize she is a different person from the original, and contracts with Kyubey to become a magical girl and save the city and remaining Pleiades.

What the rather convoluted story (made more so by the backstory being told in the order Kazumi discovers it, rather than the order in which it happened) does is literalize the unity of the witch and magical girl, which I discussed back in Chapter 9. However, by doing so it does not attempt to redeem the witch; rather, it depicts magical girls as a kind of witch. Visually, this is shown in the contrast between what we might call Kazumi's "civilian" mode of dress, characterized in the second issue as being appropriate for a much younger child, and her magical girl outfit, which is significantly skimpier than the outfits in the TV series and resembles a sexualized version of a Halloween witch costume—broad-brimmed, high-peaked black hat, tiny black skirt, bare midriff, tight blouse. Contrast Madoka's frilly pink dress, Sayaka's boyish armor, or even the most sexualized outfits in the show, Mami's high-necked, corseted dress and leggings or Kyoko's mostly unzipped sweatshirt and skirt over pants, and the distinction the comic is drawing between the innocent child and the powerful, dangerous, sexual magical girl becomes clear.

This is the greatest contrast between the comic and the series: where the unity of magical girl and witch in Madoka Magica is part of a statement about the cycle of hope and despair, in Kazumi Magica it is instead an embrace of the Madonna/whore complex. Magical girls are, in this series, all hiding terrible secrets, either that they are witches (Kazumi) or that they have done terrible things in the past (the Pleiades). Without exception, every one of the Pleiades is haunted by some misdeed from before they became magical girls, engages in deceiving Kazumi, and helps hunt and imprison other magical girls before they can become witches. This transformation is not depicted as the gradual descent into despair we saw Sayaka experience in Madoka Magica, but as simply what happens when one uses too much magic; in other words, the witch is not a product of despair but of the exercise of feminine power. Even more blatantly, the early villains are magical girls with the power to directly transform women into pseudo-witches.

The true witches, when we eventually see them, are also quite different from Madoka Magica. There, witches were implied to be mindless, delusional, and trapped within their barriers; in Kazumi Magica, they can interact with and experience the world directly, and some have personality and intellect. They are thus not Faust references and women's despair made flesh, but grotesque transformations of the magical girls' bodies. Essentially the only difference, in Kazumi Magica, between a magical girl and a witch is that the magical girl is "cute" or "sexy," while the witch is deformed or physically monstrous. The ugliness of the witch, in other words, is not her association with ugly emotions, but that she is no longer attractive to the Male Gaze, which the comic employs far more frequently than the TV series. 

The ending, too, denies the nature of the witch as depicted in Madoka Magica. Once Kazumi and her two surviving friends are again working within Kyubey's system, the comic has Kazumi's narration baldly assert that they will not become witches. The hope-despair cycle of Madoka Magica is denied entirely; instead, we have a depiction of feminine power as being permitted to continue only so long as it is constrained within a role of combating negative expressions of feminine power. As long as magical girls limit themselves to fighting witches and do not try to challenge Kyubey or his corrupt system, the comic asserts, they may continue to live and wield their power.

Ultimately, Kazumi Magica falls neatly for the trap of (hopefully unintentionally) asserting that Kyubey's system is necessary to constrain women from wielding power in ways that might challenge the status quo.
